\section{Faith Made Living by Charity}

The life judged by Christ is a life summoned and sustained by grace. Ephesians describes sinners as dead in their offenses and God as rich in mercy. He makes them alive in Christ, gives salvation through faith, and creates them anew for good works (Eph 2:1--10). Forgiveness and new life come from his initiative. Through them he enables the obedience which was impossible to our unaided strength. The good works for which we must answer are themselves a response made possible by his gift.

Paul calls us God's workmanship, created in Christ for the good works in which he intends us to walk. Grace heals our freedom for love. The dishonest person receives help to tell the truth; the resentful person receives help to forgive. That help calls for consent and action. When a familiar temptation returns, we can ask for grace and use the means of resistance available to us. Our dependence on God continues through the effort which obedience requires.

St.~James speaks with equal directness. A brother or sister lacks clothing and daily food. Someone offers peaceful words but withholds the necessities of life. James asks what such words accomplish and calls faith without works dead (Jas 2:14--26). The professed concern leaves the person hungry. James requires the believer to respond with the goods needed, when those goods are his to give. In this example the refusal concerns both a person deprived of help and the truth of the believer's profession.

Judgment according to works therefore concerns the reality of our response to God. Good intentions which never become obedience cannot indefinitely stand in place of the obedience they promise. Nor does a Christian life consist in an external record severed from faith and charity. The same act may be undertaken for vanity, fear, convenience, or the love of God. Christ knows its source and its circumstances. He judges both the act and the heart, the thing done and the freedom by which it was done.

The Church speaks of genuine merit in this life of grace. No one can merit the initial gift of forgiveness and justification. Once adopted in Christ, however, the faithful can act under the initiative of the Holy Spirit, and their graced acts have true merit before God. The merit is first attributable to grace and then to the free cooperation grace makes possible (CCC 2008--2011). God promises to reward the love which his grace enables. The response is freely ours, and we owe the ability to make it to him.

A father can supply the material for a child's gift and help the child make it. The child still gives freely. Adoption in Christ is far greater than this earthly relation, but the comparison helps explain how dependence and a real offering can coexist. We may thank God for a victory over temptation, including the strength to endure its cost. We can also ask for the reward he has promised. Gratitude acknowledges where the good began and how we were able to choose it.

Chrysostom sees this generosity in the final judgment of Matthew~25. Christ promises a kingdom in response to bread, water, clothing, shelter, and a visit. The good is real, but the gift bestowed far exceeds it. Chrysostom therefore says that those who receive the crown receive it by grace (\emph{Homilies on Matthew} LXXIX.2). The Lord receives a little service offered in charity and brings his servant into a blessedness no earthly act could purchase. This disproportion is a reason to trust his promise and offer the service within our power.

Illness can make this teaching especially consoling. A person who can undertake little outward activity may still love, pray, and receive help with gratitude. Accepting dependence can itself require an effort of humility. Limited strength changes what can reasonably be done, while leaving opportunities to consent to grace. A judgment of life's fruit based chiefly on visible activity would overlook these acts and the charity from which they proceed.

The teaching also corrects the temptation to preserve a favored religious achievement as a substitute for present conversion. We may have undertaken a great sacrifice years ago, defended the faith courageously, or served the Church with unusual generosity. These goods deserve thanksgiving. They cannot authorize a present refusal of God's command. The relationship remains living. The Lord calls the person who once obeyed to obey today, and the person who once fell to rise today. The account of a life includes perseverance, repentance, and the growth or decline of charity through time.

When a duty exceeds our strength, prayer can ask for the help to fulfill it. Having asked, we must also attempt the good and accept the assistance God gives, including help from other people. Perseverance may require repeated prayer and repeated effort after a failure. Christ knows both our dependence and our free response. His promise of reward gives us reason to continue even when the effort receives little recognition.

